% Input fragment; use the English Volume 0 preamble.
% Assign a unique label before integration; remove unused optional blocks.
\chapter{Proposed Chapter Title}\label{template:v0}
\begin{chaptersummary}
[AUTHOR TO COMPLETE: connect the established result, remaining question,
present choice, later use, and one core idea. Base the overview on the chapter.]
\end{chaptersummary}
\learninggoals{[State 2--4 assessable goals: construct, derive, compare, or verify.]}
\paragraph{Core question} [Which specific question remains unresolved?]

\section{Problem Origin and Assumptions}
[Identify the discrete operator, matrix class, variables, scales, and assumptions.]

\section{Key Relations and a Minimal Example}
[Insert the approved mathematical framework. Motivate equations and interpret results.]
[Provide a hand-checkable or reproducible example and acceptance measures.]

% Optional: retain only for a genuinely missing conceptual motivation.
\conceptbridge{[State the reader's specific question.]}
{[AUTHOR TO COMPLETE: explain the need for this step and the selection rationale.]}
{[Established object] $\longrightarrow$ [new object]}

\section{Cost and Computational Interface}
\costnote{[Define operation, storage, or communication counts, conditions, and comparison measures.]}
% Optional: \coderef{existing/relative_code_path.py}
% Optional: \crossref{A specific result needed from another chapter or volume.}
\aiexercise{[Name notation, assumptions, starting point, and a checkable calculation or example.]}
{[Change one parameter or assumption to test the mechanism.]}
% Volume 0 studyroute takes TWO arguments.
\studyroute{[Name omitted proofs, algebra, or implementation details and verified references.]}
{[Precise, established English search terms]}
